\documentclass{article}
\usepackage[utf8]{inputenc}
\usepackage{amsmath}
\usepackage{amssymb}
\usepackage{graphicx}
\usepackage{tabularx}
\usepackage{booktabs}
\usepackage{caption}
\usepackage[english]{babel}
\usepackage{setspace}
\usepackage{geometry}
\usepackage{float}
\usepackage[justification=centering]{caption}


% APA style formatting
\geometry{a4paper, margin=1in}
\setstretch{1.5}
\captionsetup[table]{labelfont=bf, justification=raggedright, singlelinecheck=false}
\captionsetup[figure]{font=small, labelfont=bf, justification=raggedright, singlelinecheck=false}

\title{Cloud Native Streaming Platform Netflix-inspired}
\author{Santiago  \\ Nicolas \\ Andrés \\ Erick Santiago Buitrago Peña   }
\date{}

\begin{document}

\maketitle

\begin{abstract}
It should be just one paragraph, with three main sentences. These sentences should answer: (i) what is the context of the problem? (ii) what do you propose
as a solution for the problem? (iii) what are the relevant results of your work?. Ensure your abstract is understandable without reading the rest of the paper. Avoid
technical jargon and acronyms unless absolutely necessary.

\end{abstract}

\section{Introduction}
This should be a section of about one page. Here, the context of
the problem should be fully described. Also, previous solutions to the same problem
should be referenced. Most of the bibliography is cited here. While referencing previous solutions, you can also mention computational techniques used, which provides a complete reference for the reader. Usually, this section does not include images or
tables.

\section{Methods and Materials}
Here you should describe the design of your solution, your
technical decisions, and why you believe your solution is appropriate. Some general
images are recommended to enhance your explanation, and this section is expected
to be 1 to 2 pages long. It is not recommended to include code here, but if you are
presenting a complex algorithm, you may add it. Remember to write full paragraphs,
with one idea per paragraph, and avoid using itemized lists. Aim for clarity and
readability. Include enough detail for reproducibility. Briefly mention any limitations
or assumptions in your approach.

\subsection{Subsection 1}
We can use different subsections depending of teacher specifications.

\begin{itemize}
    \item \textbf{Item 1}: Example 1
    \item \textbf{Item 2}: Example 2
\end{itemize}

\subsection{Subsection 2}
Subsection example 1

\subsection{Subsection 3}
Subsection example 2

\section{Results and Discussion}
Explain the experiments you performed to demonstrate
that your solution works. Summarize your unit tests in a paragraph, mentioning their
philosophy, quantity, and results. Also, discuss any integration and acceptance tests to
address software quality and user expectations. If possible, compare your results with
other solutions using metrics. Use tables to summarize test definitions and results, and
charts (such as box-plots or time series) for comparisons, depending on your project.
Highlight the most important findings. Every figure and table must have a descriptive
caption and be referenced in the text.

\section{Conclusion}
Write a couple of paragraphs to summarize your work, results, and
achievements. Clearly state why your work was successful and how it helped to solve
the problem. End with suggestions for future work or open questions that remain

\begin{thebibliography}{9}
\bibitem{rummery1994}
Rummery, G. A., \& Niranjan, M. (1994). On-line Q-learning using connectionist systems. \textit{Technical Report}, Cambridge University Engineering Department.


\end{thebibliography}

\end{document}